\documentclass[11pt]{amsart}

\usepackage[T1]{fontenc}
\usepackage{lmodern}
\usepackage{microtype}
\usepackage{amsmath,amssymb,amsthm}
\usepackage[hidelinks]{hyperref}

\hypersetup{
  pdftitle={An Exterior-Point Proof of Two Completed-Interlacing Conjectures}
}

\newtheorem{theorem}{Theorem}[section]
\newtheorem{lemma}[theorem]{Lemma}
\theoremstyle{remark}
\newtheorem{remark}[theorem]{Remark}

\newcommand{\R}{\mathbb{R}}
\newcommand{\sgn}{\operatorname{sgn}}

\title[Two completed-interlacing conjectures]
{An Exterior-Point Proof of Two Completed-Interlacing Conjectures}

\date{}

\subjclass[2020]{33C45, 42C05}
\keywords{zero interlacing, Meixner--Pollaczek polynomials,
pseudo-Jacobi polynomials}

\begin{document}

\begin{abstract}
We prove Conjectures~1 and~2 of Jordaan and Kumar on two-point completed
interlacing for Meixner--Pollaczek and pseudo-Jacobi polynomials.  The common
mechanism is an exterior-point criterion: in the relevant mixed recurrence,
the zero of the linear coefficient lies outside the interval between the two
completion points.  For pseudo-Jacobi polynomials, we also give a
three-term-recurrence proof, independent of orthogonality, of the
consecutive-degree interlacing needed throughout the conjectured parameter
range.
\end{abstract}

\maketitle

\section{Introduction}

For monic polynomials $G$ and $Q$ of degrees $n+1$ and $n$, respectively, we
write $G\prec Q$ when their zeros are real and simple and, in increasing order,
satisfy
\[
 y_1<q_1<y_2<\cdots<q_n<y_{n+1}.
\]
Jordaan and Kumar developed a general framework for completing interlacing by
two additional points~\cite{JordaanKumar}.  In the configuration relevant here,
their Theorem~2.2(e) yields the desired completed interlacing provided that,
inside the gap of the higher-degree polynomial containing a completion point,
two consecutive zeros of the lower-degree polynomial straddle that point: with
$E_1,E_2$ the two completion points, $y_{1,n+1}<\cdots<y_{n+1,n+1}$ the zeros of
the higher-degree polynomial and $z_{1,n}<\cdots<z_{n,n}$ those of the
lower-degree one, the hypothesis is
\[
 y_{k,n+1}<z_{l,n}<E_1<z_{l+1,n}<y_{k+1,n+1}
\]
for the first of the two conclusions, and the analogous condition at $E_2$ for
the second.  Their Conjectures~1 and~2 assert, for Meixner--Pollaczek and
pseudo-Jacobi polynomials, that the corresponding completed interlacing holds
without this extra placement hypothesis whenever the two completion points lie
in distinct interior gaps of the higher-degree polynomial.

We prove both conjectures.  The key observation is that each mixed recurrence
has the form
\[
 A(x)P(x)=-(x-c)G(x)+(x-r_-)(x-r_+)Q(x),
\]
and the distinguished point $c$ lies outside $[r_-,r_+]$.  A direct sign count
then forces the required local zero configuration and proves the completed
interlacing without assuming in advance that $P$ is real-rooted or coprime to
$G$.  The two conjectures are vacuous for $n=1$, so the applications below are
stated for $n\ge2$.

In the pseudo-Jacobi case the conjectured parameter range is strictly wider than
the range of the corresponding result in~\cite{JordaanKumar}: Conjecture~2
allows every $a$ satisfying \eqref{eq:PJrange}, whereas
\cite[Corollary~6.1]{JordaanKumar} assumes $a<-n-2$.  The residual band
\[
 -n-2<a<\frac{-3n-5-\sqrt{n^2+2n+5}}4
\]
is nonempty for every $n\ge1$, since $n+3>\sqrt{n^2+2n+5}$, and it is exactly
the band on which the labels $E_1,E_2$ of~\cite{JordaanKumar} occur in reverse
order; see Remark~\ref{rem:ordering}.

\section{An exterior-point criterion}

\begin{lemma}[Exterior-point criterion]\label{lem:criterion}
Let $P,Q,G$ be monic real polynomials of degrees $n,n,n+1$, respectively, with
$G\prec Q$.  Suppose
\begin{equation}\label{eq:mixed-general}
 A(x)P(x)=-(x-c)G(x)+(x-r_-)(x-r_+)Q(x),
\end{equation}
where $A(x)>0$ for every $x\in\R$ and $r_-<r_+$.  Assume that $r_-$ and $r_+$
lie in distinct interior gaps of the zeros of $G$ and that
\[
 (c-r_-)(c-r_+)>0.
\]
Then
\begin{align*}
 c<r_-&\quad\Longrightarrow\quad
 (x-r_-)G(x)\prec(x-r_+)P(x),\\
 c>r_+&\quad\Longrightarrow\quad
 (x-r_+)G(x)\prec(x-r_-)P(x).
\end{align*}
\end{lemma}

The positivity hypothesis on $A$ is less general than it appears.  Since
$G,Q,P$ are monic, the two degree-$(n+2)$ terms on the right of
\eqref{eq:mixed-general} cancel, so $\deg AP\le n+1$; a polynomial $A$ is
therefore of degree at most one, and positivity on $\R$ forces $A$ to be a
positive constant.  This is what both applications below supply.

\begin{proof}
Let $y_1<\cdots<y_{n+1}$ be the zeros of $G$, put
$R(x)=(x-r_-)(x-r_+)$, and set $F=RP$.  Since $r_\pm$ lie in gaps of $G$ and
$c\notin[r_-,r_+]$, evaluating \eqref{eq:mixed-general} at $r_\pm$ shows that
$P(r_\pm)\ne0$.  At each $y_j$ we have
\begin{equation}\label{eq:Fsign}
 A(y_j)F(y_j)=R(y_j)^2Q(y_j)\ne0.
\end{equation}
Because $G\prec Q$, the values $Q(y_j)$ alternate in sign.  Hence $F$ has a
zero in every interval $(y_j,y_{j+1})$.

Suppose first that $c<r_-$, and let $r_-\in(y_k,y_{k+1})$.  The gap containing
$r_+$ lies to the right, so
\[
 R(y_k)>0>R(y_{k+1}).
\]
Monicity and $G\prec Q$ give
\[
 \sgn Q(y_k)=\sgn G(r_-)=-\sgn Q(y_{k+1}).
\]
Using \eqref{eq:mixed-general} at $y_k,y_{k+1}$ and at $r_-$ therefore yields
\[
 \sgn P(y_k)=\sgn P(y_{k+1})=\sgn G(r_-),
 \qquad
 \sgn P(r_-)=-\sgn G(r_-).
\]
Thus $P$ has a zero in each of $(y_k,r_-)$ and $(r_-,y_{k+1})$.  Consequently,
$F$ has at least three distinct zeros in this gap and at least one in each of
the other $n-1$ interior gaps.  Since $\deg F=n+2$, these are all its zeros and
all are simple.  In particular, the gap containing $r_+$ has $r_+$ as its
unique zero of $F$.  Removing the zero $r_-$ from $F=(x-r_-)(x-r_+)P$ shows
that the zeros of $(x-r_+)P$ occupy exactly the intervals between consecutive
zeros of $(x-r_-)G$.  This proves the first conclusion.  The case $c>r_+$ is
symmetric.
\end{proof}

\section{Meixner--Pollaczek polynomials}

Let $P_m^{(\lambda)}(x;\phi)$ be the monic Meixner--Pollaczek polynomial,
normalised by $P_{-1}^{(\lambda)}=0$, $P_0^{(\lambda)}=1$ and the
three-term recurrence
\begin{equation}\label{eq:MPrec}
 P_{m+1}^{(\lambda)}(x;\phi)
 =\bigl(x+(m+\lambda)\cot\phi\bigr)P_m^{(\lambda)}(x;\phi)
 -\frac{m(m+2\lambda-1)}{4\sin^2\phi}\,P_{m-1}^{(\lambda)}(x;\phi),
\end{equation}
valid for $\lambda>0$ and $\phi\in(0,\pi)$.  This fixes the normalisation used
below; for the family and its orthogonality weight see
\cite[\href{https://dlmf.nist.gov/18.19}{\S18.19}]{DLMF}.
For $n\ge2$ and $\lambda>0$, set
\[
 \phi_c=\frac12\arccos\!\left(\frac{2\lambda-n-1}{2\lambda+n+1}\right)
 \in\left(0,\frac\pi2\right)
\]
and assume
\[
 \phi\in(0,\phi_c)\cup(\pi-\phi_c,\pi).
\]
The quadratic
\begin{equation}\label{eq:T}
 T(x)=(1-\cos2\phi)x^2+(n+1)\sin2\phi\,x
 +\lambda^2-\bigl(\lambda^2+(n+1)\lambda\bigr)\cos2\phi
\end{equation}
has two distinct real roots $E_1<E_2$ by
\cite[Lemma~5.1]{JordaanKumar}.  Here $T$ is~\cite[(5.3)]{JordaanKumar} and the
displayed range of $\phi$ is~\cite[(5.4)]{JordaanKumar}, while the explicit
roots \cite[(5.7)--(5.8)]{JordaanKumar} are labelled there with $E_1<E_2$, so
the two labellings agree.  Conjecture~1 of~\cite{JordaanKumar} reads as
follows, quoted in the source's notation.

\begin{quote}
\emph{Let $\{y_{k,n+1}\}_{k=1}^{n+1}$ denote the zeros
of $P_{n+1}^{(\lambda+1)}(x;\phi)$.  Let $E_1, E_2$ be as defined
in~\textup{(5.7)} and~\textup{(5.8)}, respectively, for $\lambda>0$ and $\phi$
satisfying~\textup{(5.4)}.  If there exist distinct $k, k' \in \{1, \dots, n\}$
such that $y_{k,n+1} < E_1 < y_{k+1,n+1}$ and
$y_{k',n+1} < E_2 < y_{k'+1,n+1}$, then
\begin{gather*}
 (x-E_1)P_{n+1}^{(\lambda+1)}(x;\phi) \prec (x-E_2)P_{n}^{(\lambda)}(x;\phi)\\
 \text{or}\quad
 (x-E_2)P_{n+1}^{(\lambda+1)}(x;\phi) \prec (x-E_1)P_{n}^{(\lambda)}(x;\phi)
\end{gather*}
holds.}
\end{quote}

\noindent
Its hypothesis says that $E_1$ and $E_2$ lie in distinct interior gaps of the
zeros of $P_{n+1}^{(\lambda+1)}(x;\phi)$, and its conclusion is the disjunction
of the two alternatives below.

\begin{theorem}\label{thm:MP}
Suppose $E_1$ and $E_2$ lie in distinct interior gaps of the zeros of
$P_{n+1}^{(\lambda+1)}(x;\phi)$.  Then
\begin{align*}
 0<\phi<\phi_c
 &\quad\Longrightarrow\quad
 (x-E_2)P_{n+1}^{(\lambda+1)}(x;\phi)
 \prec (x-E_1)P_n^{(\lambda)}(x;\phi),\\
 \pi-\phi_c<\phi<\pi
 &\quad\Longrightarrow\quad
 (x-E_1)P_{n+1}^{(\lambda+1)}(x;\phi)
 \prec (x-E_2)P_n^{(\lambda)}(x;\phi).
\end{align*}
In particular, Conjecture~1 of~\cite{JordaanKumar} holds.
\end{theorem}

\begin{proof}
Write
\[
 P=P_n^{(\lambda)}(x;\phi),\qquad
 G=P_{n+1}^{(\lambda+1)}(x;\phi),\qquad
 Q=P_n^{(\lambda+1)}(x;\phi),\qquad
 c=\lambda\cot\phi.
\]
The mixed recurrence in~\cite[(5.10)]{JordaanKumar}, together with
$1-\cos2\phi=2\sin^2\phi$, is
\begin{equation}\label{eq:MPmixed}
 \frac{(2\lambda+n)(2\lambda+n+1)}{2(1-\cos2\phi)}P
 =-(x-c)G+(x-E_1)(x-E_2)Q.
\end{equation}
Both sides of \eqref{eq:MPmixed} have degree at most $n+1$, the
degree-$(n+2)$ terms on the right cancelling as noted after
Lemma~\ref{lem:criterion}; the identity may therefore be checked from
\eqref{eq:MPrec} at $n+2$ points.

The left coefficient is positive.  Also $G$ and $Q$ are Meixner--Pollaczek
polynomials of consecutive degrees with the same parameters, hence orthogonal
with respect to a common weight on $\R$
\cite[\href{https://dlmf.nist.gov/18.19}{\S18.19}]{DLMF}, and the zeros of
consecutive orthogonal polynomials separate each other
\cite[\href{https://dlmf.nist.gov/18.2.vi}{\S18.2(vi)}]{DLMF}; thus
$G\prec Q$.
Since $T(x)=(1-\cos2\phi)(x-E_1)(x-E_2)$, direct substitution gives
\[
 (c-E_1)(c-E_2)
 =\frac{T(c)}{1-\cos2\phi}
 =\frac{\lambda(2\lambda+n+1)}{1-\cos2\phi}>0.
\]
Thus Lemma~\ref{lem:criterion} applies.  Moreover,
\[
 T'(c)=(2\lambda+n+1)\sin2\phi.
\]
Because $T$ opens upward and $c\notin[E_1,E_2]$, the derivative is positive
exactly when $c>E_2$ and negative exactly when $c<E_1$.  It is positive on
$(0,\phi_c)$ and negative on $(\pi-\phi_c,\pi)$, so the two conclusions follow.
\end{proof}

\section{Pseudo-Jacobi polynomials}

Let $P_m(x;\alpha,b)$ denote the monic pseudo-Jacobi polynomial; see
\cite{JordaanTookos,JordaanKumar}.  As in the previous section we fix the
normalisation by a recurrence: throughout, $P_m(x;\alpha,b)$ is determined by
$P_{-1}=0$, $P_0=1$ and \eqref{eq:PJrec} below.  We first record a proof of the
consecutive-degree interlacing needed later, using nothing beyond that
recurrence.

\begin{lemma}\label{lem:PJinterlace}
Let
\[
 n\ge1,\qquad b\in\R,\qquad
 \alpha<-n-\tfrac12,\qquad \alpha\ne-n-1.
\]
Then
\[
 P_{n+1}(x;\alpha,b)\prec P_n(x;\alpha,b).
\]
\end{lemma}

\begin{proof}
We first record that \eqref{eq:PJrec} is the recurrence of the classical
family, so that the normalisation fixed above is the usual one; the interlacing
argument itself uses \eqref{eq:PJrec}--\eqref{eq:gamma} alone.
Comparing the pseudo-Jacobi definition with the Jacobi hypergeometric formula
\cite[\href{https://dlmf.nist.gov/18.5.E7}{(18.5.7)}]{DLMF} gives
\[
 P_m(x;\alpha,b)
 =i^{-m}\widehat P_m^{(\alpha+ib,\alpha-ib)}(ix),
\]
where $\widehat P_m^{(\mu,\nu)}$ is the monic Jacobi polynomial.  The monic
Jacobi recurrence coefficients
\cite[\href{https://dlmf.nist.gov/3.5.E33_1}{(3.5.33\_1)}]{DLMF} are rational in
the parameters and therefore extend wherever their denominators are nonzero.
Substitution yields, for $0\le m\le n$,
\begin{equation}\label{eq:PJrec}
 P_{m+1}(x;\alpha,b)
 =\left(x+\frac{\alpha b}{(m+\alpha)(m+\alpha+1)}\right)P_m(x;\alpha,b)
 -\gamma_mP_{m-1}(x;\alpha,b),
\end{equation}
where $P_{-1}=0$, $P_0=1$, and
\begin{equation}\label{eq:gamma}
 \gamma_m
 =-\frac{m(m+2\alpha)\bigl((m+\alpha)^2+b^2\bigr)}
 {(m+\alpha)^2(2m+2\alpha-1)(2m+2\alpha+1)}.
\end{equation}
The stated assumptions make every denominator in
\eqref{eq:PJrec}--\eqref{eq:gamma} nonzero.  For $1\le m\le n$, the three
factors $m+2\alpha$, $2m+2\alpha-1$, and $2m+2\alpha+1$ are negative, and hence
$\gamma_m>0$.

Thus $P_m$ is the characteristic polynomial of the leading $m\times m$
principal block of a real symmetric tridiagonal matrix whose off-diagonal
entries are $\sqrt{\gamma_m}$.  The matrix is irreducible, so strict
principal-minor interlacing gives the result.  Equivalently, equality at a
putative common zero would propagate backward through \eqref{eq:PJrec} to
$P_0=0$, a contradiction.
\end{proof}

Assume henceforth that
\begin{equation}\label{eq:PJrange}
 a<\frac{-3n-5-\sqrt{n^2+2n+5}}4,
 \qquad a\ne-n-2,
\end{equation}
and
\begin{equation}\label{eq:brange}
 |b|>
 \sqrt{\frac{4(a+n+1)^2(a+n+2)^2(-2a-n-2)}
 {(n+1)\left|4(a+n+1)^2-(2n-2)(a+n+1)-(n+1)\right|}}.
\end{equation}
Define
\begin{equation}\label{eq:R}
 R_{n,a,b}(x)
 =x^2-\frac{b(n+1)}{(a+n+1)(a+n+2)}x
 +\frac{(a+n+1)(a+n+2)(2a+n+2)+b^2(n+1)}
 {(a+n+1)(a+n+2)(2a+2n+3)}.
\end{equation}
Under \eqref{eq:PJrange}--\eqref{eq:brange}, this polynomial has two distinct
real roots by~\cite[Lemma~6.1]{JordaanKumar}; write them in increasing order as
$e_-<e_+$.  Here \eqref{eq:R}, \eqref{eq:PJrange} and \eqref{eq:brange}
are~\cite[(6.2), (6.3), (6.4)]{JordaanKumar}, and the explicit roots
\cite[(6.8)--(6.9)]{JordaanKumar} are labelled there $E_1,E_2$.  Conjecture~2
of~\cite{JordaanKumar} reads as follows, quoted in the source's notation.

\begin{quote}
\emph{Let $\{y_{k,n+1}\}_{k=1}^{n+1}$ be the zeros of
$P_{n+1}(x;a+1,b)$. Let $E_1$ and $E_2$ be given by~\textup{(6.8)}
and~\textup{(6.9)}, where $a$ and $b$ satisfy the corresponding
conditions~\textup{(6.3)} and~\textup{(6.4)}.  If there exist distinct
$k, k' \in \{1, \dots, n\}$ such that $y_{k,n+1} < E_1 < y_{k+1,n+1}$ and
$y_{k',n+1} < E_2 < y_{k'+1,n+1}$, then
\begin{gather*}
 (x-E_1) P_{n+1}(x;a+1,b) \prec (x-E_2)P_{n}(x;a,b)\\
 \text{or}\quad
 (x-E_2)P_{n+1}(x;a+1,b) \prec (x-E_1)P_{n}(x;a,b)
\end{gather*}
holds.}
\end{quote}

\noindent
Again the hypothesis says that the two completion points lie in distinct
interior gaps of the zeros of $P_{n+1}(x;a+1,b)$, and the conclusion is a
disjunction; see Remark~\ref{rem:ordering} for the labelling.

\begin{theorem}\label{thm:PJ}
Let $n\ge2$, and suppose $e_-$ and $e_+$ lie in distinct interior gaps of the
zeros of $P_{n+1}(x;a+1,b)$.  Set
\[
 \Delta_c=\frac{b(2a+n+3)}{(a+n+1)(a+n+2)}.
\]
Then $\Delta_c\ne0$ and
\begin{align*}
 \Delta_c<0
 &\quad\Longrightarrow\quad
 (x-e_-)P_{n+1}(x;a+1,b)
 \prec (x-e_+)P_n(x;a,b),\\
 \Delta_c>0
 &\quad\Longrightarrow\quad
 (x-e_+)P_{n+1}(x;a+1,b)
 \prec (x-e_-)P_n(x;a,b).
\end{align*}
In particular, Conjecture~2 of~\cite{JordaanKumar} holds.
\end{theorem}

\begin{proof}
Put
\[
 u=a+n+1,\qquad c=\frac bu,
\]
and
\[
 P=P_n(x;a,b),\qquad
 G=P_{n+1}(x;a+1,b),\qquad
 Q=P_n(x;a+1,b).
\]
The mixed recurrence~\cite[(6.10)]{JordaanKumar}, valid whenever its
coefficients are defined, is
\begin{equation}\label{eq:PJmixed}
 AP=-(x-c)G+R_{n,a,b}(x)Q,
\end{equation}
where
\[
 A=\frac{2(2u-n-1)(2u-n)(u^2+b^2)}
 {u(2u-1)(2u)(2u+1)}.
\]
The bound \eqref{eq:PJrange} is strictly smaller than $-n-\tfrac32$, so
$u<-\tfrac12$ and therefore $A>0$.  It also gives
$a+1<-n-\tfrac12$, while $a\ne-n-2$ gives $a+1\ne-n-1$.
Lemma~\ref{lem:PJinterlace} therefore yields $G\prec Q$.

A direct calculation gives
\begin{equation}\label{eq:Rc}
 R_{n,a,b}(c)
 =\frac{(2u-n)(u^2+b^2)}{u^2(2u+1)}>0.
\end{equation}
Since $R_{n,a,b}$ is monic with roots $e_-<e_+$, it follows that
$c\notin[e_-,e_+]$, and Lemma~\ref{lem:criterion} applies.  Finally,
\[
 R_{n,a,b}'(c)=\Delta_c.
\]
For an upward-opening quadratic at a point outside its root interval,
$R'(c)<0$ exactly when $c<e_-$ and $R'(c)>0$ exactly when $c>e_+$.
The stated alternatives follow, and $\Delta_c\ne0$.
\end{proof}

\begin{remark}[Ordering of the completion points]\label{rem:ordering}
The formulas labelled $E_1,E_2$ in~\cite[(6.8)--(6.9)]{JordaanKumar} are
in increasing order when $a<-n-2$ and in reverse order on the remaining
admissible band.  Using $e_-<e_+$ removes this notational reversal.
Conjecture~2 is a disjunction of the two displayed interlacings and is
invariant under exchanging the two labels, so Theorem~\ref{thm:PJ} proves the
conjecture exactly as stated.
\end{remark}

\small
\begin{thebibliography}{9}

\bibitem{DLMF}
NIST Digital Library of Mathematical Functions,
\url{https://dlmf.nist.gov/}, Release 1.2.7 of 2026-06-15,
F.~W.~J. Olver, A.~B. Olde Daalhuis, D.~W. Lozier, B.~I. Schneider,
R.~F. Boisvert, C.~W. Clark, B.~R. Miller, B.~V. Saunders, H.~S. Cohl,
and M.~A. McClain, eds.

\bibitem{JordaanKumar}
K.~Jordaan and V.~Kumar,
\emph{Two-point completion of zero interlacing and Wronskian-type bounds,
with applications to Jacobi and other orthogonal polynomials},
\href{https://arxiv.org/abs/2604.25692v2}
{arXiv:2604.25692v2 [math.CA]}, 2026.

\bibitem{JordaanTookos}
K.~Jordaan and F.~To\'okos,
\emph{Orthogonality and asymptotics of pseudo-Jacobi polynomials for
non-classical parameters},
J. Approx. Theory \textbf{178} (2014), 1--12,
\href{https://doi.org/10.1016/j.jat.2013.10.003}
{doi:10.1016/j.jat.2013.10.003}.

\end{thebibliography}

\end{document}